\chapter{اللقاء الخامس عشر: تابع تفسير سورة البقرة (٤)}

\section{تأويل (ولما جاءهم كتاب من عند الله)}
السلام عليكم ورحمة الله وبركاته. الحمد لله رب العالمين، والصلاة والسلام على نبينا محمد وعلى آله وصحبه وسلم. هذا هو الدرس الثالث عشر من اجتماعنا على قراءة ودراسة تفسير الإمام \rawi{ابن جرير الطبري} عليه رحمة الله.

وصلنا بحمد الله إلى قول الله عز وجل: \begin{ayah}وَلَمَّا جَاءَهُمْ كِتَابٌ مِّنْ عِندِ اللَّهِ مُصَدِّقٌ لِّمَا مَعَهُمْ وَكَانُوا مِن قَبْلُ يَسْتَفْتِحُونَ عَلَى الَّذِينَ كَفَرُوا فَلَمَّا جَاءَهُم مَّا عَرَفُوا كَفَرُوا بِهِ ۚ فَلَعْنَةُ اللَّهِ عَلَى الْكَافِرِينَ\end{ayah}.

\isnad{قال الإمام ابن جرير الطبري رحمه الله:} «يعني جل ثناؤه بقوله: \textit{ولما جاءهم كتاب من عند الله}، ولما جاء اليهود من بني إسرائيل كتاب من عند الله، يعني بالكتاب القرآن الذي أنزله على محمد ﷺ. \textit{مصدق لما معهم}، يعني مصدق للذي معهم من الكتب التي أنزلها الله من قبل القرآن».

ثم ساق \rawi{الطبري} الأخبار عن \rawi{قتادة} و\rawi{الربيع} في أن المعنى هو القرآن الذي أنزل على محمد ﷺ مصدق لما معهم من التوراة والإنجيل.

\section{تأويل (وكانوا من قبل يستفتحون على الذين كفروا)}
\isnad{قال أبو جعفر:} «أي وكان اليهود... يستفتحون بمحمد ﷺ، ومعنى الاستفتاح الاستنصار، ويستنصرون الله به على مشركي العرب من قبل مبعثه».
وروى \rawi{الطبري} بإسناده عن \rawi{عاصم بن عمرو بن قتادة} وعن \rawi{ابن عباس} أن اليهود كانوا يستفتحون على الأوس والخزرج برسول الله ﷺ قبل مبعثه، فلما بعثه الله من العرب كفروا به وجحدوا ما كانوا يقولون فيه.

\begin{fawaid}[أهمية السياق الزمني والمكاني في التفسير]
من المهم أن نقرأ الحدث في سياقه الزمني وبيئته الجغرافية؛ ففي ذلك الوقت كان اليهود هم أهل الكتاب والعلم، وكانوا يستطيلون على مشركي العرب ويستفتحون عليهم ببعثة النبي الخاتم، فلما جاء النبي ﷺ بالقرآن ووافق بعض ما عندهم وخالف بعضه وفضح أسرارهم، كانت هذه أبلغ حجة عليهم وعلى المشركين في إثبات صدق نبوته ﷺ.
\end{fawaid}

\subsection{الخلاف النحوي في جواب (لما)}
طرح \rawi{الطبري} تساؤلاً لغوياً: «فأين جواب قوله: ولما جاءهم كتاب من عند الله؟ قيل: قد اختلف أهل العربية في جوابه».
فذكر قول \rawi{الأخفش} أنه مما تُرِك جوابه استغناء بمعرفة المخاطبين به، كما في قوله: \begin{ayah}وَلَوْ أَنَّ قُرْآنًا سُيِّرَتْ بِهِ الْجِبَالُ\end{ayah}. وذكر قول \rawi{الفراء} أن الجواب في الفاء التي في قوله: \textit{فلما جاءهم ما عرفوا كفروا به}، كقولك: لما قمت فلما جئتنا أحسنت.

\begin{fawaid}[معرفة مصادر الأقوال اللغوية]
معرفة مصادر العلم ومنشأ الأقوال أصل في دراسة أي علم. ويهمنا جداً في كتاب الطبري أن نعرف من الذين رجع إليهم في الأمور اللغوية أو النحوية، كالاعتماد على كتب "مجاز القرآن" لأبي عبيدة، و"معاني القرآن" للأخفش، و"معاني القرآن" للفراء. فالطبري يذكر أقوالهم وقد يعزو إليهم أو لا يعزو، لكن معرفة هذه المصادر تعين الطالب على وضع الأقوال في سياقها وفهم مرادها.
\end{fawaid}

\separator

\section{تأويل (فلما جاءهم ما عرفوا كفروا به فلعنة الله على الكافرين)}
\isnad{قال أبو جعفر:} «البيان الواضح أنهم تعمدوا الكفر بمحمد ﷺ بعد قيام الحجة بنبوته عليهم، وقطع الله عذرهم بأنه رسوله إليهم... فخزي الله وإبعاده على الجاحدين ما قد عرفوا من الحق».

\begin{fawaid}[الفرق بين العلم والمعرفة في الاستعمال القرآني]
المعرفة هي مطلق الإدراك، ولا يلزم أن يعمل بها صاحبها، كما قال تعالى: \textit{يعرفونه كما يعرفون أبناءهم}، \textit{يعرفون نعمة الله ثم ينكرونها}. أما العلم في القرآن فهو الإدراك المتبوع بالخشية والعمل والاتباع، كما قال: \textit{إنما يخشى الله من عباده العلماء}. فهؤلاء اليهود عرفوا الحق وأدركوه، لكنهم لم يعملوا به فكفروا.
\end{fawaid}

\separator

\section{تأويل (بئسما اشتروا به أنفسهم أن يكفروا بما أنزل الله بغيا)}
\isnad{قال أبو جعفر:} «ومعنى قوله جل ثناؤه بئسما اشتروا به أنفسهم: ساء ما اشتروا به أنفسهم».
ثم ذكر اختلاف أهل العربية في إعراب (بئسما)، فرجح أن (بئسما) مرفوع بالراجع من الهاء في قوله (اشتروا به)، وأن (أن يكفروا) مترجمة عن (بئسما)، فيكون المعنى: «بئس الشيء باع اليهود به أنفسهم: كفرهم بما أنزل الله بغياً وحسداً أن ينزل الله من فضله».

وأوضح أن معنى \textit{اشتروا} في هذا الموضع: باعوا، فالعرب تستعمل الشراء بمعنى البيع والعكس. و\textit{بغياً} يعني: تعدياً وحسداً للنبي ﷺ وللعرب أن الله جعل النبوة فيهم.

\begin{fawaid}[أسباب الضلال وأمانة العلم]
الضلال لا يكون دائماً بسبب الجهل، بل قد يكون سببه الحسد أو الكبر أو حب الدنيا، كما حصل مع اليهود الذين كفروا بالنبي ﷺ حسداً من عند أنفسهم. وهذا يعلم طالب العلم أن العلم أمانة وميثاق ومسؤولية من الله، فلا ينبغي أن يستعمله لنيل حظوظ النفس من مال أو شهرة، أو أن يبغي به على الناس ويستعلي عليهم، فمن استعمل علمه في غير طاعة الله ونفع عباده كان كمن اشترى الضلالة بالهدى.
\end{fawaid}

\begin{fawaid}[التفريق بين التمحيص والإهانة]
في تفسير قوله تعالى: \textit{وللكافرين عذاب مهين}، بيّن الطبري أن العذاب المهين هو الذي يورث صاحبه ذلة وهواناً ويخلد فيه، وهو مخصص للكفار. أما ما يصيب المؤمنين من عذاب أو نكال (كإقامة الحدود في الدنيا أو عذاب العصاة في الآخرة قبل دخول الجنة) فهو للتمحيص والتطهير وليس للإهانة، ثم يوردهم معدن العز والكرامة.
\end{fawaid}

\separator

\section{تأويل (وإذا قيل لهم آمنوا بما أنزل الله قالوا نؤمن بما أنزل علينا)}
\isnad{قال أبو جعفر:} «وإذا قيل لليهود... صدقوا بما أنزل الله من القرآن على محمد ﷺ، قالوا نصدق بالتوراة التي أنزل الله على موسى. \textit{ويكفرون بما وراءه} أي ويجحدون بما سوى التوراة وما بعدها من كتب الله... \textit{وهو الحق مصدقا لما معهم} أي أن القرآن حق موافق لما في أيديهم من التوراة في الأمر باتباع محمد ﷺ».

\subsection{تأويل (قل فلم تقتلون أنبياء الله من قبل إن كنتم مؤمنين)}
\isnad{قال أبو جعفر:} «قل يا محمد لليهود: لم تقتلون -إن كنتم يا معشر اليهود مؤمنين بما أنزل الله عليكم- أنبياءه؟ وقد حرم الله في الكتاب الذي أنزل عليكم قتلهم... وذلك من الله جل ثناؤه تكذيب لهم في قولهم: نؤمن بما أنزل علينا».

\begin{fawaid}[دقة الحجاج القرآني]
هذه الآية من أبلغ الحجج التي لقنها الله لنبيه ﷺ ليرد بها على اليهود. فهم ادعوا أنهم يؤمنون بما أنزل عليهم ويكتفون به، فرد الله عليهم بأن إيمانهم بكتابهم يقتضي اتباع الأنبياء وتوقيرهم، فكيف تدعون الإيمان وأنتم تقتلون أنبياء الله الذين ابتعثوا إليكم؟ فبان كذبهم في دعوى الإيمان بكتابهم فضلاً عن غيره.
\end{fawaid}

\separator

\section{تأويل (ولقد جاءكم موسى بالبينات ثم اتخذتم العجل من بعده وأنتم ظالمون)}
\isnad{قال أبو جعفر:} «أي جاءكم بالبينات الدالة على صدقه... ثم اتخذتم العجل إلهاً من بعد أن فارقكم موسى... وأنتم واضعون للعبادة في غير موضعها، لأن العبادة لا تنبغي لغير الله».

\begin{fawaid}[الآيات لا توجب الإيمان للمكابرين]
هذه الآية وما شابهها تُبين أن مجيء الآيات والبينات ليس موجباً للإيمان إلا لمن أراد الهدى وقصد الحق. فبنو إسرائيل رأوا عصا موسى وانفلاق البحر وغير ذلك من البينات العظام، ومع ذلك بمجرد أن غاب عنهم موسى عبدوا العجل! فهذا تسلية وتوجيه للنبي ﷺ وللدعاة بأن كثرة الحجج والآيات لا تكفي وحدها لهداية القلوب المعاندة.
\end{fawaid}

\separator

\section{تأويل (واشربوا في قلوبهم العجل بكفرهم)}
اختلفوا في معنى إشراب العجل في القلوب؛ فرجح \rawi{الطبري} أنه: \textbf{أُشربوا في قلوبهم حُبَّ العجل}، وأن الحذف تم للعلم به، كما يقال: (واسأل القرية) أي أهلها. وذكر قولاً آخر للسدي: أن موسى سحل العجل وذراه في الماء وأمرهم أن يشربوا منه، فمن كان يحبه خرج على شاربه الذهب.
واختار \rawi{الطبري} التأويل الأول لكونه المعنى المتبادر في لغة العرب، إذ يقال "أُشرب قلب فلان حب كذا" إذا خالطه وغلب عليه.

\separator

\section{تأويل (قل إن كانت لكم الدار الآخرة عند الله خالصة... فتمنوا الموت)}
\isnad{قال أبو جعفر:} «هذه الآية مما احتج الله به لنبيه محمد ﷺ على اليهود... وفضح بها أحبارهم... وقال لفريق اليهود: إن كنتم محقين فتمنوا الموت، فإن ذلك غير ضاركم إن كنتم محقين فيما تدعون من الإيمان وقرب المنزلة من الله، بل إن أُعطيتم أمنيتكم من الموت صرتم إلى الراحة... فامتنعت اليهود من إجابة النبي ﷺ إلى ذلك لعلمها أنها إن تمنت الموت هلكت... \textit{ولن يتمنوه أبداً بما قدمت أيديهم} أي بما كسبت أيديهم من الذنوب والكفر وتكذيب محمد ﷺ».

\separator

\section{تأويل (ولتجدنهم أحرص الناس على حياة ومن الذين أشركوا)}
\isnad{قال أبو جعفر:} «يقول يا محمد لتجدن أشد الناس حرصاً على الحياة في الدنيا وأشدهم كراهة للموت: اليهود... وإنما وصفهم الله بذلك لعلمهم بما أُعد لهم في الآخرة على كفرهم... \textit{ومن الذين أشركوا} أي وهم أحرص من الذين أشركوا الذين لا يؤمنون بالبعث».

\subsection{تأويل (يود أحدهم لو يعمر ألف سنة)}
\isnad{قال الطبري:} «هذا خبر من الله عن الذين أشركوا... يود أحد هؤلاء... لو يعمر في الدنيا ألف سنة، حتى جعل بعضهم تحية بعض: عش ألف عام... \textit{وما هو بمزحزحه من العذاب أن يعمر} أي وما طول العمر بمبعده ولا منجيه من عذاب الله».

\separator

\section{تأويل (قل من كان عدوا لجبريل فإنه نزله على قلبك بإذن الله)}
أورد \rawi{الطبري} الإجماع على أن الآية نزلت جواباً لليهود حين زعموا أن جبريل عدو لهم وأن ميكائيل ولي لهم. وذكر سببين لنزولها: الأول مناظرة جرت بينهم وبين النبي ﷺ، والثاني مناظرة جرت بينهم وبين عمر بن الخطاب رضي الله عنه.

\isnad{قال أبو جعفر:} «فأعلمهم الله أن من كان لجبريل عدواً فإن الله له عدو، وأنه من الكافرين، فنص عليه باسمه وعلى ميكائيل باسمه، لئلا يقول قائل: لسنا لله ولا لملائكته باعداء... ليقطع بذلك تلبيسهم».

\separator

\section{تأويل (أوكلما عاهدوا عهدا نبذه فريق منهم)}
رد \rawi{الطبري} على من زعم أن الواو والفاء في (أوكلما) و (أفكلما) زائدتان لا معنى لهما، \isnad{فقال:} «وقد بيّنا فيما مضى أنه غير جائز أن يكون في كتاب الله حرف لا معنى له». فرجح أنها واو عطف أُدخلت عليها ألف الاستفهام.

\begin{fawaid}[لا يوجد حرف زائد بلا معنى في القرآن]
من القواعد العظيمة في التفسير والتي يقررها الطبري ويرد بها على اللغويين المعتزلة: أنه لا يجوز أن يوجد في كتاب الله المعجز حرف زائد أو حشو لا معنى له. فكل حرف وُضع لمعنى مقصود، والقول بالزيادة المطلقة طعن في بلاغة القرآن وحكمته، وقد أطال ابن تيمية رحمه الله في تقرير هذا الأصل وبيان أن ما يسميه النحويون حروف "صلة" إنما تفيد التوكيد أو نفي المجاز ونحو ذلك.
\end{fawaid}

\separator

\section{تأويل (ولما جاءهم رسول من عند الله... نبذ فريق... كتاب الله وراء ظهورهم)}
\isnad{قال أبو جعفر:} «نبذ فريق، يعني جحدوه ورفضوه... وجعلوه وراء ظهورهم، وهذا مثل يقال لكل رافض أمراً... \textit{كأنهم لا يعلمون}، وهذا من الله أخبار عنهم أنهم جحدوا الحق على علم منهم به، وعاندوا أمر الله فخالفوه».

\separator

\section{تأويل (واتبعوا ما تتلوا الشياطين على ملك سليمان)}
\isnad{قال أبو جعفر:} «واتبعوا، يعني اليهود الذين نبذوا كتاب الله... ما تتلو الشياطين من السحر، على ملك سليمان، أي في عهد ملك سليمان. \textit{وما كفر سليمان} بنسبة السحر إليه، \textit{ولكن الشياطين كفروا يعلمون الناس السحر}».

\subsection{تأويل (وما أنزل على الملكين ببابل هاروت وماروت)}
اختلف المفسرون في (ما) هل هي نافية للجحد (بمعنى: ولم ينزل على الملكين)، أم هي موصولة (بمعنى: الذي أنزل)؟
رجّح \rawi{الطبري} بقوة أن (ما) بمعنى (الذي)، ورد على من جعلها نافية مبيناً فساد زعمهم من وجوه لغوية ومعنوية. وذكر أن هاروت وماروت ملكان أُهبطا ابتلاءً وفتنة للناس، يعلمان السحر بعد التحذير منه \textit{وما يعلمان من أحد حتى يقولا إنما نحن فتنة فلا تكفر}.

وأورد \rawi{الطبري} في ذلك آثاراً وروايات كثيرة بعضها من الإسرائيليات في شأن الملكين وكيف فُتنا. ولم يجزم الطبري بحقيقة السحر هل هو تخييل ومخاريق أم له حقيقة تقلب الأعيان، وإنما أورد الخلاف في ذلك.

\begin{fawaid}[هل يصح أن تُعلم الملائكة السحر؟]
أجاب الطبري عن إشكال كيف تعلم الملائكة السحر وهو تفريق وإضرار؟ بأن الله يعرّف عباده الخير والشر، ويمتحنهم بما يشاء. فتعليم الملكين للسحر كان بأمر الله لهما ابتلاءً وفتنة للناس، وكانا يطيعان الله بامتثال أمره في ذلك وبتحذيرهما للناس (إنما نحن فتنة)، ويمحص الله المؤمنين بترك التعلم، ويخزي الكافرين باتباع السحر.
\end{fawaid}

\separator

بهذا ننهي لقاءنا، وغداً إن شاء الله نكمل باقي الآيات من سورة البقرة. والسلام عليكم ورحمة الله وبركاته.